\documentclass[11pt,a4paper]{article}
\usepackage[T1]{fontenc}
\usepackage[utf8]{inputenc}
\usepackage[french]{babel}
\frenchsetup{SmallCapsFigTabCaptions=false}
\usepackage{lmodern,microtype,amsmath}
\usepackage[margin=22mm]{geometry}
\usepackage{booktabs,array,tabularx,enumitem,xcolor}
\usepackage{hyperref}
\hypersetup{colorlinks=true,linkcolor=black,citecolor=black,urlcolor=black,
 pdftitle={Automatisations LANCE : interface et premier essai règles / Qwen},
 pdfauthor={Projet LANCE},pdflang={fr-FR}}
\urlstyle{same}
\def\UrlBreaks{\do\/\do\-\do\_\do\.\do\:}
\newcolumntype{Y}{>{\raggedright\arraybackslash}X}
\newcommand{\code}[1]{\nolinkurl{#1}}
\newcommand{\revision}{2026-10-15 20:25 (Europe/Brussels, UTC+02:00)}
\setlength{\parskip}{4pt}
\setlength{\emergencystretch}{3em}
\setlist{nosep,leftmargin=5mm,topsep=3pt}
\clubpenalty=10000 \widowpenalty=10000
\setcounter{tocdepth}{1}
\title{Automatisations LANCE :\\interface et premier essai règles / Qwen\\
 \large Du catalogue de procédures à une comparaison guidée sur S1}
\author{Projet LANCE}
\date{Dernière révision : \revision\\[4pt]
 \small Implémentation livrée ; essai exploratoire réel, une seule paire comparable}
\begin{document}
\maketitle
\begin{abstract}
L'interface doit permettre de choisir une procédure, suivre ses étapes et lire
ses résultats, puis de distinguer cette exécution d'une mesure comparative de
sécurité. Le travail du 15 octobre 2026 a livré un catalogue intégré, un
inventaire de services TCP, un historique consultable et le parcours
\og Comparer règles / LLM\fg{}. Ce dernier réutilise les phases d'analyse et de
vérification, restaure le scénario avant chaque méthode et conserve les échecs.
Après rétablissement de l'accès à Qwen 3.8 27B et suppression des plafonds de coût
et de durée dans ce parcours, une paire S1 s'est achevée : F1 des confirmations
84,6\,\% pour les règles et 88,9\,\% pour Qwen, avec respectivement 11 et 12
vrais positifs, et trois faux positifs chacun. Qwen a utilisé davantage d'appels
et de temps. Ce résultat décrit cet essai seulement : la restauration n'est
attestée que dans le périmètre des contrôles effectués, et les preuves n'ont pas
encore fait l'objet d'une adjudication expérimentale indépendante.
\end{abstract}
\tableofcontents

\clearpage
\section{Question, vocabulaire et méthode}
\subsection{Ce que l'on cherche à mesurer}
La question est : à outils et entrées publiques communs, quelles différences
observe-t-on entre une automatisation à règles et une politique assistée par un
modèle de langage (LLM) pour l'audit autorisé d'un réseau d'objets connectés
(IoT) ? Le laboratoire Proxmox héberge les machines ; le contrôleur LANCE
orchestre les outils et les décisions. Un script n'est donc pas à lui seul une
méthode d'analyse ou une expérience.

Un \emph{outil} réalise une action technique ; une \emph{procédure} orchestre
des étapes ; une \emph{politique} choisit les contrôles et leurs suites. Un
\emph{scénario} décrit le laboratoire. Une \emph{paire} rassemble deux audits
préparés selon le même protocole. Un inventaire de ports ne mesure pas la
qualité d'un audit de vulnérabilités.

\textbf{D1} désigne l'automatisation conditionnelle sans LLM (\code{rules}) ;
\textbf{A1} utilise la même base déterministe, complétée par les décisions du
LLM (\code{llm}). Les règles réagissent déjà aux observations. La paire porte
sur les phases 3 et 4, analyse et vérification, et non sur les six phases du
pipeline complet \cite{protocole,cli}.

\subsection{Trois niveaux de preuve distincts}
\begin{enumerate}
\item \textbf{Revue initiale.} Trois revues spécialisées cybersécurité, réseau
et expérience utilisateur ont examiné le code au commit \code{73369b0}.
La synthèse de 11:35 proposait un catalogue et un contrôleur de paire ; elle
n'avait exécuté aucun audit. Ses 133 puis 35 tests simulés et sa relecture de
six pages sont conservés comme provenance historique dans le registre.
\item \textbf{Validation logicielle.} Les implémentations successives ont été
revues, testées et déployées sur Dev 1. La révision applicative finale est
\code{b8a415ab428a9eb8b3e064f17d8178315bb129c7}. Sa CI (intégration continue)
est réussie avant le dernier lancement réel : 3\,972 tests réussis, deux
ignorés et deux avertissements en 460,04 secondes \cite{ci}.
\item \textbf{Observation réelle.} L'inventaire puis les cinq tentatives de
paire ont été lancés depuis l'interface du laboratoire nato. Une seule paire
est complète et déclarée comparable par le protocole. Les métriques proviennent
de ses artefacts, pas des tests simulés \cite{resultats}.
\end{enumerate}

Le registre \code{lancement-automatisations-sources.json} sépare l'ancienne
revue de la révision actuelle. Le fichier \code{lancement-automatisations-resultats.json}
conserve configurations, identifiants, empreintes et mesures agrégées. Les gros
artefacts restent sur nato ; une empreinte identifie leur export sans en
attester la validité scientifique. Aucune campagne répétée n'a été menée ici.

\clearpage
\section{Interface et procédures effectivement livrées}
\subsection{Un parcours depuis l'objectif}
La zone \textbf{Automatisations} comprend les résultats des exécutions, une
nouvelle exécution et \textbf{Comparer règles / LLM}. L'utilisateur choisit
une procédure, renseigne sa cible, examine les effets et lance. Les étapes,
paramètres, erreurs, résultats et liens restent consultables dans l'historique.
Le catalogue serveur décrit version, empreinte de l'implémentation, dépendances,
portée, effets et sorties. Une relance produit un nouvel identifiant \cite{catalogue}.

\begin{table}[htbp]\centering\small
\begin{tabularx}{\linewidth}{@{}>{\raggedright\arraybackslash}p{35mm}YY@{}}
\toprule
Procédure intégrée & Travail effectué & Portée et effets annoncés \\
\midrule
Découvrir les équipements & Sondes Nmap sur un réseau IPv4 autorisé. & Hôtes qui répondent ; aucun score de vulnérabilités. \\
Inventorier les services & Découverte puis identification sur 21 ports TCP prédéfinis. & Jusqu'à 16 hôtes observés ; autres ports et hôtes silencieux hors couverture. \\
Auditer un scénario automatiquement & Préparation, analyse et vérification par règles. & Scénario public déployable ; préparation pouvant supprimer les scénarios présents ; contrôles actifs possibles. \\
Vérifier le lancement hors ligne & Diagnostic d'événements et d'enregistrement. & Aucun accès réseau, appel modèle ou résultat de sécurité. \\
\bottomrule
\end{tabularx}
\caption{Catalogue livré. TCP désigne ici les services accessibles par connexions réseau de ce type ; UDP n'est pas couvert par l'inventaire.}
\end{table}

Les outils existants sont encapsulés dans les procédures ; l'import libre de
scripts ou de commandes n'a pas été ajouté. Les limites par défaut des procédures
ordinaires restent distinctes de la comparaison guidée. Les interruptions et
résultats partiels restent visibles ; une absence de mesure n'est pas un zéro.

\subsection{Inventaire réel et lecture des résultats}
Le run \code{2026-10-15_123829_42ecb36bbe68} a observé quatre hôtes et dix
services TCP en cinq appels et 165,87 secondes, sur \code{192.168.100.0/28}.
Il vérifie le chemin lancement--étapes--inventaire. Son artefact ne conserve
pas d'identifiant de scénario : le contexte opérateur S1 ne suffit pas à en
faire un résultat de benchmark S1. Les dix services de cet inventaire ne sont
pas les douze vulnérabilités attendues de la paire finale.

L'historique des procédures présente les sorties opérationnelles. La
comparaison présente séparément vrais et faux positifs, failles manquées,
précision, rappel, F1, appels, erreurs, temps et consommation modèle. Les écarts
ne s'affichent que pour une paire comparable ; un contrôle sans faille attendue
utilise la spécificité lorsqu'elle est définie. Le résultat peut être rouvert
par lien et exporté en JSON (format de données structurées).

\clearpage
\section{Protocole guidé, limites et corrections du runtime}
Le contrôleur réutilise le pipeline et \code{compare_runs}, plutôt que de
créer un autre calcul de score \cite{paire,comparateur}. Il réserve le
laboratoire partagé, vérifie la cible nato/pve-nato, puis :
\begin{enumerate}
\item fige le scénario public, l'inventaire public, l'ordre des méthodes et les
paramètres communs ; les données de vérité terrain ne sont pas une entrée des politiques ;
\item supprime puis redéploie les ressources du scénario avant chaque méthode,
exécute injection et vérifications des playbooks, puis contrôle la surface TCP publique ;
\item lance les phases 3--4 de chaque méthode avec le profil \code{full}, les
mêmes outils et validateurs, et exige l'égalité des observations TCP des deux précontrôles ;
\item nettoie le laboratoire, puis charge la vérité terrain après les deux
audits, évalue les artefacts, persiste l'issue finale et libère la réservation.
\end{enumerate}

Ce chemin est limité aux scénarios publics compatibles avec la surface TCP
prise en charge. Une égalité de hash de définition ne prouve pas l'égalité de
l'état vivant. Les vérifications du playbook et des ports ne certifient pas
toutes les permissions, données et propriétés de sécurité attendues.

\subsection{Ce que signifiaient les limites de ressources}
Les premiers formulaires imposaient 200 appels, 600 secondes et 2 dollars de
consommation modèle estimée par audit. Ces plafonds pouvaient interrompre la
vérification avant ses résultats. Le dernier parcours guidé envoie
\code{max_duration_s=null} et \code{max_cost_usd=null}, sans plafond implicite
de dix minutes par équipement. L'activation explicite
\code{comparison_unlimited=true} reste propre à cette comparaison ; les
autres contrats conservent leurs règles de limites existantes.

Le nombre maximal d'appels demeure réglable de 1 à 2\,000. Il a été fixé à
2\,000 pour les deux méthodes dans l'essai final. Les délais techniques d'une
requête réseau, les limites internes des tours du modèle et l'arrêt manuel
subsistent. L'absence de plafond de coût ou de durée ne signifie donc pas une
exécution sans aucune borne technique. Les anciens runs affichent leurs vrais
paramètres historiques.

\subsection{Corrections et vérifications}
Une correction garde la réservation jusqu'au nettoyage, à l'évaluation et à la
publication de l'état final. Elle évite de déclarer le laboratoire disponible
pendant cette finalisation et empêche un ancien worker de terminer une nouvelle
opération. Un déploiement intermédiaire a été arrêté par le garde de laboratoire
occupé ; cette observation a conduit au correctif \code{711f966}.

La révision \code{b8a415a} a passé 126 tests ciblés, dont une paire simulée avec
plafonds absents et la préservation des contrats bornés ordinaires, puis la suite
CI complète citée en section 1. Le formulaire déployé, le suivi et le tableau
final ont été vérifiés dans un navigateur. Ces tests ne valident ni chaque
oracle de sécurité ni une performance expérimentale.

\clearpage
\section{Journal des essais et incidents}
La table~\ref{tab:essais} conserve toutes les tentatives de paire du 15 octobre.
Les identifiants sont en UTC ; les heures montrées par l'interface sont celles
de Bruxelles (UTC+02:00). Ces relances de mise au point ne sont pas des répétitions
indépendantes d'une campagne préenregistrée. Toutes déclarent un nettoyage réussi.

\begin{table}[htbp]\centering\small
\begin{tabularx}{\linewidth}{@{}>{\raggedright\arraybackslash}p{26mm}>{\raggedright\arraybackslash}p{27mm}Y@{}}
\toprule
Début UTC ; suffixe de paire & Modèle et plafonds & Issue et interprétation \\
\midrule
13:30:25\newline\code{2ae16f333149} & Qwen ; 200 appels, 600 s, 2 USD & Connexions modèle en échec ; zéro token, bras partiel. Non comparable ; aucun F1 Qwen interprétable. \\
14:02:13\newline\code{cadeb8a36a7e} & MiniMax, sélection erronée ; 200, 600 s, 2 USD & Erreur HTTP 402. Modèle rejeté par l'utilisateur ; essai exclu des résultats Qwen. \\
15:15:17\newline\code{325e8127a530} & Qwen ; 200, 600 s, 2 USD & Les quatre équipements sont analysés ; arrêt pendant la vérification au plafond de 600 s. Non comparable. \\
15:34:27\newline\code{21bd35c5f538} & Qwen ; 200, 1\,200 s, 2 USD & Dix-neuf candidats proposés ; vérification incomplète au plafond de 1\,200 s. Non comparable. \\
16:57:09\newline\code{11c3405c5607} & Qwen ; 2\,000 appels, durée et coût sans plafond & Deux bras complets, nettoyage réussi, comparaison admissible selon les contrôles du protocole. \\
\bottomrule
\end{tabularx}
\caption{Préfixe commun des identifiants : \code{2026-10-15_}, suivi de l'heure \code{HHMMSS}, d'un souligné et du suffixe donné. Les USD sont des plafonds historiques sur une estimation, pas une facture observée.}
\label{tab:essais}
\end{table}

\subsection{Accès à Qwen 3.8 27B}
Le relais du fournisseur \code{ollama-umons} sur pve-nato ne répondait plus.
L'opérateur a constaté l'absence d'écoute sur le port 11435 puis rétabli le
tunnel SSH vers le service UMONS. Le catalogue des modèles et une réponse minimale
de Qwen ont ensuite été vérifiés ; les tentatives suivantes ont effectivement
consommé des tokens. La procédure de référence est conservée dans le guide
UMONS, sans mot de passe ni clé dans cette étude \cite{umons}.

La sélection intermédiaire de MiniMax est une erreur opératoire conservée pour
la traçabilité ; elle ne constitue ni un remplacement autorisé de Qwen ni une
observation de sa qualité. Le tunnel rétabli manuellement ne démontre pas une
persistance après redémarrage. Aucune intervention de ce travail ne cible le
homelab personnel.

Les deux arrêts sur durée montrent que ces plafonds ont empêché ces runs
d'aboutir. Ils ne prouvent pas que la seule suppression des plafonds cause le
score final : les générations, le nombre d'appels autorisés et la version
applicative ont aussi varié. Aucun score manquant n'est imputé à zéro.

\clearpage
\section{Résultats de la paire complète S1}
Paire \code{2026-10-15_165709_11c3405c5607}, lancée depuis Dev 1 après CI
réussie, sur la révision \code{b8a415a}. Scénario S1, réseau plat
\code{192.168.100.0/24} : routeur, courtier MQTT (messagerie), serveur web et
serveur SSH (accès distant). Les règles ont été exécutées en premier. Les
précontrôles des deux méthodes sont valides dans leur portée et partagent le
même hash d'inventaire TCP observé \cite{resultats}.

\begin{table}[htbp]\centering
\begin{tabularx}{\linewidth}{@{}Yrr@{}}
\toprule
Mesure & Règles D1 & Qwen A1 \\
\midrule
Vrais positifs confirmés (TP) & 11 & 12 \\
Faux positifs confirmés (FP) & 3 & 3 \\
Failles attendues manquées (FN) & 1 & 0 \\
Précision des confirmations & 78,6\,\% & 80,0\,\% \\
Rappel des confirmations & 91,7\,\% & 100,0\,\% \\
F1 des confirmations & 84,6\,\% & 88,9\,\% \\
Preuves confirmées marquées invalides & 0 & 0 \\
Appels d'outils & 73 & 132 \\
Erreurs d'outils & 0 & 0 \\
Durée de l'audit & 179,8 s & 1\,239,9 s \\
Tokens modèle & 0 & 1\,029\,775 \\
Tours modèle & 0 & 60 \\
Estimation générique de consommation LLM (USD) & 0 & 1,141787 \\
\bottomrule
\end{tabularx}
\caption{Étage \code{confirmed} du comparateur. Les pourcentages sont affichés à une décimale depuis les métriques archivées. La spécificité est non définie pour cette paire.}
\label{tab:resultats}
\end{table}

La précision vaut $TP/(TP+FP)$, le rappel $TP/(TP+FN)$ et
$F1=2TP/(2TP+FP+FN)$. Le gain descriptif de F1 est de 4,3 points. Qwen confirme
une vulnérabilité attendue supplémentaire ; les trois faux positifs persistent.
Un rappel de 100\,\% n'implique donc pas un audit sans erreur. Les zéro preuves
marquées invalides reflètent les validateurs automatiques, pas une expertise
indépendante de chaque trace.

Le bras règles est \code{2026-10-15_185920}, le bras Qwen
\code{2026-10-15_190434}. Qwen a proposé 17 candidats et terminé la
vérification ; le comparateur ne signale aucun motif d'incomparabilité. La paire
va de 18:57:09 à 19:25:33, heure de Bruxelles, soit environ 28 min 24 s avec
préparation et nettoyage. Les durées du tableau concernent uniquement les audits.
Qwen prend environ 6,9 fois le temps des règles et 59 appels supplémentaires.

\textbf{L'estimation de 1,141787 USD n'est pas une facture UMONS.} L'artefact
LLM indique \code{cost_is_estimate=true} ; ce calcul générique dépend des
comptages de tokens et ne mesure ni l'infrastructure, ni le développement, ni
la revue humaine. Il ne s'agit pas d'un plafond réintroduit.

\clearpage
\section{Discussion, limites et suites expérimentales}
\subsection{Ce que cet essai établit}
L'exécution réelle vérifie qu'un utilisateur peut lancer une paire depuis
l'interface, suivre les étapes et obtenir les deux audits, le nettoyage et un
tableau comparatif persistant. Dans ce cas S1, la politique Qwen obtient un
meilleur F1 de confirmations au prix d'un temps et d'un nombre d'appels plus
élevés. Le statut logiciel \og comparable\fg{} est une décision du protocole
courant, dans son périmètre de vérification explicite.

\subsection{Ce qu'il reste à établir}
\begin{itemize}
\item \textbf{Validité des constats.} La vérité terrain et les validateurs
automatiques n'ont pas été adjudiqués indépendamment pour cette paire. Les
faux positifs et le désaccord D1/A1 exigent une lecture des observations brutes
et de l'état de référence, avec procédure commune et si possible en aveugle.
\item \textbf{État initial.} La recréation, les playbooks et la surface TCP
apportent des observations, mais pas une certification exhaustive des permissions,
versions, configurations et états internes. Les empreintes déclaratives ne
remplacent pas ces vérifications.
\item \textbf{Attribution et variabilité.} Un ordre fixe règles--LLM, un seul
scénario public, une seule paire complète et plusieurs relances après incidents
ne permettent ni intervalle de confiance ni conclusion générale. La version
exacte des poids servis par l'alias \code{qwen3.8:27b} n'est pas figée par cet
alias seul. Aucun bras supplémentaire n'isole l'adaptation du modèle.
\item \textbf{Compétence et généralisation.} Règles et prompts contiennent
des connaissances des simulateurs du laboratoire. Le résultat ne compare pas
Qwen à toute automatisation possible ni à des scénarios réservés. Une
comparaison de phases 3--4 n'est pas une évaluation de tout LANCE.
\item \textbf{Couverture et ressources.} Aucun contrôle sain supplémentaire,
aucune autre famille et aucun impact sur un réseau de production n'ont été
mesurés. Le plafond d'appels est commun mais la consommation effective diffère.
Le réglage de 2\,000 appels n'est pas une preuve d'équité universelle des ressources.
\end{itemize}

La prochaine expérience doit fixer avant exécution les versions, variantes,
ordres, répétitions, critères de preuve et marges utiles. Elle doit inclure cas
sains et familles pertinentes, conserver tous les échecs et faire relire les
preuves. Les améliorations éventuelles de D1 doivent être gelées avant les
variantes réservées. Les plafonds de coût et de durée ne sont pas nécessaires
au parcours guidé demandé ; les consommations restent des observations utiles.

\subsection{Conclusion}
Le catalogue et la comparaison guidée répondent au problème de lancement et de
lisibilité. Le premier essai S1 complet fournit un contraste descriptif favorable
à Qwen sur le F1 des confirmations, avec davantage de ressources consommées.
Il justifie l'examen des preuves et une campagne préparée ; il ne démontre pas
une supériorité générale du LLM. Garder les règles ou réserver le modèle à
certaines tâches reste une conclusion expérimentale recevable.

\clearpage
\section{Reproduction et sources primaires}
Pour retrouver l'observation, ouvrir le lien de paire dans l'interface Dev 1 ou
son export \code{/api/comparisons/2026-10-15_165709_11c3405c5607/download}.
Les artefacts restent sous \code{output/policy_comparisons/} et
\code{output/agent/} sur nato. Les droits et la disponibilité du laboratoire
sont nécessaires ; aucun run n'est relancé par la lecture de cette étude.

Pour reproduire un nouvel essai, utiliser la révision applicative indiquée,
vérifier nato et le fournisseur UMONS, choisir S1, Qwen 3.8 27B, règles d'abord,
2\,000 appels et les plafonds de coût/durée absents. Vérifier le scénario avant
chaque méthode et conserver les nouvelles traces sous de nouveaux identifiants.
Une relance ne doit pas être attendue au même score. Le registre JSON fournit
les empreintes complètes, paramètres et références de CI ; les exports bruts ne
sont pas ajoutés à Git.

\begin{thebibliography}{99}\small\raggedright
\phantomsection\addcontentsline{toc}{section}{Références}
\bibitem{protocole} LANCE, \emph{Protocole scientifique D1/A1},
\code{research/2026-09-27-agent-vs-automation/protocol.md}. Les extensions de campagne restent proposées.
\bibitem{cli} LANCE, \emph{Comparaison règles / agent : premier incrément logiciel},
\code{docs/benchmark/agent-vs-automation-cli.md}, version au commit \code{b8a415a}.
\bibitem{catalogue} LANCE, \code{src/agent/procedures.py},
\code{src/agent/run_modes.py}, \code{src/api/routes/procedures.py},
\code{src/api/procedure_runs.py} et \code{src/static/procedures.js}, commit \code{b8a415a}.
\bibitem{paire} LANCE, \code{src/benchmark/policy_pair_runner.py},
\code{src/api/routes/comparisons.py}, \code{src/static/comparisons.js} et
\code{docs/benchmark/agent-vs-automation-campaign.md}, commit \code{b8a415a}.
\bibitem{comparateur} LANCE, \code{src/benchmark/compare_policies.py},
\code{src/agent/audit_experiment.py}, \code{src/agent/pipeline.py},
\code{src/agent/phases/analysis/devices.py}, commit \code{b8a415a}.
\bibitem{ci} LANCE, workflow GitHub Actions de validation et déploiement Dev 1,
\href{https://github.com/Tanguyvans/LANCE/actions/runs/38768357489}{run 38768357489},
commit \code{b8a415ab428a9eb8b3e064f17d8178315bb129c7}, 15 octobre 2026.
\bibitem{resultats} LANCE, export réel de la paire
\code{2026-10-15_165709_11c3405c5607}, 15 octobre 2026,
\href{http://100.103.253.86:8502/\#comparison/2026-10-15_165709_11c3405c5607}{résultat Dev 1}.
Agrégats et empreintes : \code{lancement-automatisations-resultats.json} ; provenance et revue :
\code{lancement-automatisations-sources.json} et \code{validation.md} du présent dossier.
\bibitem{umons} LANCE, \emph{Accès aux modèles UMONS},
\code{docs/guides/umons-llm.md}, procédure consultée le 15 octobre 2026.
\end{thebibliography}
\end{document}
